\chapter*{Notation}

The following notational conventions are used throughout these notes.

\begin{center}
\begin{tabular}{ll}
    \textbf{Notation} & \textbf{Meaning} \\
    \hline
    $a$, $x$, $\lambda$ & scalars (italic letters) \\
    $\mathbf{v}$, $\mathbf{x}$ & vectors (bold lowercase letters); $v_i$ is the $i$-th entry of $\mathbf{v}$ \\
    $\mathbf{A}$, $\mathbf{B}$ & matrices (bold uppercase letters); $a_{ij}$ is the entry in row $i$, column $j$ \\
    $\mathsf{x}$, $\mathsf{y}$ & random variables (sans-serif letters); realizations are italic ($x$, $y$) \\
    $\boldsymbol{\mathsf{x}}$ & random vectors (bold sans-serif letters) \\
    $\mathbb{R}$, $\mathbb{C}$ & the sets of real and complex numbers \\
    $\mathbb{R}^n$, $\mathbb{R}^{m \times n}$ & the sets of $n$-dimensional real vectors and $m \times n$ real matrices \\
    $\mathcal{N}$, $\mathcal{U}$ & probability distributions (calligraphic letters) \\
    $\mathbb{E}[\cdot]$, $\Pr(\cdot)$ & expectation and probability \\
    $\mathbf{A}^\top$ & transpose of $\mathbf{A}$ \\
    $\overline{z}$ & complex conjugate of $z$ \\
    $O(\cdot)$ & asymptotic computational complexity (big-O) \\
    \texttt{typewriter} & MATLAB code \\
\end{tabular}
\end{center}

When the same letter appears in several forms, the forms are related: for example, $x$ may denote a scalar, $\mathbf{x}$ a vector with entries $x_i$, and $\mathsf{x}$ a random variable whose realization is $x$.
